% ATIVIDADES.tex --- Os cantos do homem-sombra (1ª série do Ensino Médio)
% Paratexto pedagógico para o estudante: atividades de aprofundamento e reflexão
% (TR SEDUC-GO 118787/2026, itens 6.4.2 e 6.4.3)

\blankAtodd
\part{Atividades de aprofundamento e reflexão}

\chapter{Para começar a conversa}

Você acaba de ler uma história muito antiga, contada e recontada por
gerações de Hupd'äh, um povo que vive na Floresta Amazônica, na região
do Alto Rio Negro. Ela é curta: cabe em poucas páginas e foi publicada
lado a lado em duas línguas, o hup e o português. Mas não se engane com
o tamanho. Essa história explica a origem de algo que, para os Hupd'äh,
é muito importante até hoje: os \textit{caapivaiá}, os cantos que
animam as festas.

Ela também foi escolhida para abrir um percurso de três anos. Na
1ª série, você lê \textit{Os cantos do homem-sombra}; na 2ª,
\textit{A árvore dos cantos}, dos Yanomami; na 3ª, \textit{A terra uma
só}, escrita pelo Guarani Mbya Timóteo Verá Tupã Popygua. São três
povos diferentes, com línguas, territórios e histórias diferentes. Não
existe ``o índio'' em geral: existem mais de trezentos povos indígenas
no Brasil, e cada um tem seus modos de contar o mundo.

As atividades a seguir foram organizadas em etapas:

\begin{itemize}
\item {\formular\textbf{Antes da leitura}}: para você levantar o que já
sabe (e o que imagina saber) sobre os povos indígenas e conhecer os
Hupd'äh.
\item {\formular\textbf{Durante a leitura}}: um roteiro para ler, e
reler, com atenção.
\item {\formular\textbf{Depois da leitura}}: questões de compreensão e
interpretação.
\item {\formular\textbf{Aprofundamento e reflexão}}: textos e perguntas
que ampliam a discussão sobre oralidade, tradução, línguas indígenas e
preconceito.
\item {\formular\textbf{Criação e protagonismo}}: propostas para você
ilustrar, recontar, entrevistar, pesquisar e produzir.
\item {\formular\textbf{Pesquisa e debate}}, {\formular\textbf{Conexões
com Goiás}} e {\formular\textbf{Projeto integrador}}: para levar o que
aprendeu para a escola, para a sua família e para a sua região.
\end{itemize}

Mantenha um \textit{diário de leitura}: um caderno ou arquivo digital
onde você vai anotar impressões, dúvidas, palavras novas, desenhos e
respostas. Ele será útil em várias atividades.

\chapter{Antes da leitura}

\section{o que eu sei, o que eu imagino}

\begin{enumerate}[label=\arabic*.]
\item Em dupla, escrevam em uma folha as cinco primeiras palavras que
vêm à cabeça quando vocês ouvem a expressão ``povos indígenas''. Depois,
comparem com as listas das outras duplas e organizem, no quadro, as
palavras mais repetidas.
\item Observem a lista da turma e discutam:
\begin{enumerate}[label=\alph*)]
\item Quais palavras falam de pessoas reais, que vivem hoje, e quais
falam de um passado distante?
\item Quais palavras se referem a um lugar só (por exemplo, a floresta)?
Existem indígenas vivendo em cidades, no Cerrado, no litoral, no Sul do
país?
\item Alguma palavra carrega preconceito ou generalização? De onde vocês
acham que vêm essas ideias: da escola, da televisão, das redes sociais,
de conversas de família?
\end{enumerate}
\item Guarde a lista. Ao final do estudo do livro, vocês vão voltar a ela
para ver o que mudou.
\end{enumerate}

\section{quem são os hupd'äh}

Leia com atenção os textos iniciais do livro: ``Nota dos
organizadores'', ``Como foi feito este livro'' e ``Para ler as palavras
hup''. Depois, responda no diário de leitura:

\begin{enumerate}[label=\arabic*.,resume]
\item Onde vivem os Hupd'äh? Quantas pessoas são e em quantas aldeias
vivem, segundo a nota?
\item Como os Hupd'äh garantem sua alimentação? Que atividades cabem às
mulheres e aos homens, e em que momentos elas se misturam?
\item A nota diz que, nas viagens pela mata, os Hupd'äh encontram
``outros tipos de gente, como a \textit{gente-sombra}, a
\textit{gente-onça} e a \textit{gente-árvore}''. Antes de ler a
história, levante hipóteses: o que pode ser uma ``gente'' que não é
humana?
\item Localize a região do Alto Rio Negro em um mapa do Brasil. Use,
por exemplo, o site Terras Indígenas no Brasil
(\url{https://terrasindigenas.org.br}) ou a enciclopédia Povos
Indígenas no Brasil, do Instituto Socioambiental
(\url{https://pib.socioambiental.org}). Observe: de que país a região
faz fronteira? Que grandes rios passam por ali? Qual é a distância
aproximada entre essa região e a sua cidade?
\end{enumerate}

\section{um título, uma imagem, uma língua}

\begin{enumerate}[label=\arabic*.,resume]
\item O título em hup é \textit{B’atɨ́b’ yám pɨnɨ̀g}. Observe as letras:
quais delas não existem no alfabeto do português? Consulte ``Para ler as
palavras hup'' e descubra o que indicam o til, os acentos e o apóstrofo.
\item A nota informa que o hup é uma \textit{língua tonal}. Pesquise o
que isso significa. Em português, a entonação muda o sentido de uma
frase (compare ``Ele chegou.'' e ``Ele chegou?''). Em uma língua tonal,
o que muda quando o tom de uma sílaba muda?
\item Folheie o livro sem ler o texto, só observando as ilustrações de
Anita Ekman. Que cores predominam? Que seres, plantas e paisagens
aparecem? Com base nas imagens e no título, escreva em três linhas o que
você imagina que vai acontecer na história.
\end{enumerate}

\chapter{Durante a leitura}

Leia a história pelo menos duas vezes: a primeira, de uma vez só, para
conhecer o enredo; a segunda, devagar, observando os pontos a seguir.
Se possível, faça também uma leitura em voz alta, com colegas, dividindo
os trechos. Esta é uma história que nasceu para ser falada e ouvida.

\begin{enumerate}[label=\arabic*.]
\item {\formular\textbf{Os cantos de abertura.}} O livro começa com dois
cantos: ``Um canto do homem-sombra'', em português, e ``Way Naku'', em
hup. Leia os dois em voz alta, mesmo sem saber pronunciar tudo. Que
palavras e versos se repetem? Que efeito a repetição produz?
\item {\formular\textbf{O texto em itálico.}} Antes da história começar,
há um parágrafo em itálico que apresenta a gente-sombra. Anote no diário
as características desses seres: o que os torna perigosos e o que os
torna admiráveis.
\item {\formular\textbf{O encontro.}} Onde está o homem-sombra quando a
história começa? O que o homem hup está fazendo quando ouve a música?
Anote o lugar, o momento do dia e as ações de cada personagem.
\item {\formular\textbf{O tempo.}} Marque no texto todas as indicações
de tempo (``por muito tempo'', ``a noite toda'', ``quando o sol começou
a nascer''\ldots). Quanto tempo dura, aproximadamente, a história?
\item {\formular\textbf{Os sons.}} Circule ou copie todas as
onomatopeias (palavras que imitam sons). Depois, procure as palavras
correspondentes no texto em hup, na mesma página.
\item {\formular\textbf{A voz de quem conta.}} Em que momentos o
narrador fala de si mesmo e de seu povo? Copie esses trechos.
\item {\formular\textbf{O final.}} O que acontece com o homem-sombra? E
com o homem hup? O que muda para o povo hup depois desse encontro?
\item {\formular\textbf{Palavras novas.}} Liste as palavras que você não
conhecia (por exemplo, \textit{cucura}, \textit{igarapé},
\textit{terçado}, \textit{umari}) e procure seu sentido no Glossário ou
em um dicionário.
\end{enumerate}

\chapter{Depois da leitura: compreensão e interpretação}

Responda no diário de leitura. Sempre que possível, justifique suas
respostas com trechos do livro.

\begin{enumerate}[label=\arabic*.]
\item Releia o parágrafo em itálico que antecede a história:
\begin{quote}
Homens e mulheres-sombra são muito fortes e perigosos. [\ldots] Mas
muitos deles são sábios e conhecem cantos, mitos e benzimentos.
\end{quote}
Que palavra, nesse trecho, marca a mudança de uma ideia para a outra?
Por que essa dupla característica (perigo e sabedoria) é importante para
entender o que acontece na história?

\item No início, lemos: ``No alto de uma árvore de cucura, um
homem-sombra chamado Way Naku cantava um \textit{caapivaiá}''. Por que é
significativo que o homem-sombra tenha um nome próprio, enquanto o
homem hup é chamado apenas de ``um homem hup''?

\item ``Encantado com a música, um homem hup aproximou-se, sentou-se ao
pé da árvore e ficou escutando as canções do homem-sombra.'' O homem hup
sabia que a gente-sombra é perigosa. O que o leva a se aproximar? Que
sentimento a palavra \textit{encantado} expressa?

\item O narrador diz: ``Muitos dos cantos que nós, Hupd'äh, cantamos nas
festas foram cantados pelo Way Naku''. Quais cantos são citados? O que
os nomes desses cantos revelam sobre o ambiente e a vida dos Hupd'äh?

\item Releia: ``O homem hup ficou lá a noite toda e, depois, o outro dia
inteiro. Escutou e prendeu os cantos do homem-sombra.''
\begin{enumerate}[label=\alph*)]
\item Que sentido tem o verbo \textit{prender} nesse trecho? Substitua-o
por outro verbo e compare o efeito.
\item O que esse trecho revela sobre a forma de aprender entre os
Hupd'äh?
\end{enumerate}

\item ``Quando o sol começou a nascer, o homem hup bateu com o terçado
no tronco da árvore de cucura.'' A história não explica por que o homem
hup faz isso. Levante duas hipóteses diferentes para essa ação e
justifique cada uma com elementos do texto.

\item Compare as onomatopeias do texto em português com as do texto em
hup:

\begin{center}
\begin{tabular}{p{.35\linewidth}p{.35\linewidth}}
\textit{Português} & \textit{Hup}\\
\hline
Tac! & Täk!\\
Puffff! & Pë'!\\
Vrummm! & Mmmm'!\\
tchibum! & \textit{tapuh}!\\
\end{tabular}
\end{center}

\begin{enumerate}[label=\alph*)]
\item Alguns pares são parecidos, outros são bem diferentes. O que isso
mostra sobre a maneira como cada língua ``ouve'' e representa os sons?
\item Por que um tradutor precisa decidir se mantém ou recria uma
onomatopeia?
\end{enumerate}

\item Qual é o papel do jacaré na história? Sem ele, o desfecho seria o
mesmo? Explique.

\item Como o homem-sombra confunde o jacaré com o homem hup, o leitor
fica sabendo algo sobre a maneira como Way Naku percebe o mundo. O que
você conclui sobre ele?

\item ``Esse foi o primeiro homem hup que soube cantar os
\textit{caapivaiá}, os cantos das festas.'' Por que podemos dizer que
esta é uma \textit{narrativa de origem}? Origem de quê?

\item Releia o final:
\begin{quote}
Ele ouviu o homem-sombra cantar, aprendeu bem e começou a ensinar para
seus filhos, irmãos, cunhados. E até hoje todos cantam e dançam os
\textit{caapivaiá} que animam nossas festas.

Foi isso o que ouvi de meus avós.
\end{quote}
\begin{enumerate}[label=\alph*)]
\item Que caminho o conhecimento percorre, do homem-sombra até os
Hupd'äh de hoje?
\item Qual é a função da última frase? Por que o narrador faz questão de
dizer de quem ouviu a história?
\end{enumerate}

\item A história é narrada em terceira pessoa (``um homem hup'', ``o Way
Naku''), mas em alguns momentos aparece a primeira pessoa (``nós,
Hupd'äh'', ``nossas festas'', ``meus avós''). Que efeito essa mudança
produz? Quem é esse narrador?

\item Compare o ``Um canto do homem-sombra'' (em português) com a
história. Que ação da narrativa aparece também no canto? Por que os dois
últimos versos (``\textit{Way Naku, yari nóóy mah / Marika, Way
Naku}'') foram mantidos em hup, sem tradução?

\item Volte às hipóteses que você escreveu antes da leitura, a partir do
título e das ilustrações. Elas se confirmaram? O que mais o surpreendeu
na história?
\end{enumerate}

\chapter{Aprofundamento e reflexão}

\section{da voz dos avós ao livro: oralidade, escrita e tradução}

A história que você leu não foi inventada por um escritor sozinho,
diante de uma folha em branco. Ela foi contada por muitas gerações de
Hupd'äh antes de chegar a este livro. Em 2002, a linguista Patience Epps,
que morava na aldeia de Barreira Alta para estudar a língua hup, pediu
ao senhor Mário Andrade Pires, ``um sábio ancião hup'', que lhe contasse
a história dos cantos do homem-sombra. A narração foi gravada, escrita
e traduzida para o português com a ajuda dos professores hupd'äh. O
canto \textit{Way Naku} foi cantado em 2010, na aldeia de Tat Dëh, pelo
senhor Ponciano Salustiano Ramos para o antropólogo Danilo Paiva Ramos.

Antes de virar este livro, a narrativa fez parte de um livro na própria
língua, \textit{Hupd’äh nɨh Pɨnɨ̀gd’äh} (\textit{Contos dos Hupd'äh}),
usado nas escolas das aldeias para ensinar crianças e adultos a ler e a
escrever em hup. Ou seja: o mesmo texto circula para dentro, fortalecendo
a língua e a memória do povo, e para fora, apresentando essa cultura a
leitores não indígenas.

Quando uma história oral é escrita, algo se ganha e algo se transforma.
Ganha-se permanência: o texto pode ser relido, estudado, levado para
longe. Mas perdem-se a voz, o ritmo, os gestos, as pausas, a melodia dos
cantos e a presença de quem conta e de quem ouve. Na tradução, surgem
novas escolhas: como passar para o português um canto, um nome, uma
onomatopeia? O que fica sem tradução, e por quê?

\begin{enumerate}[label=\arabic*.]
\item Quem são os autores deste livro? Faça uma lista de todas as
pessoas que participaram de sua criação (narradores, cantores,
organizadores, tradutores, ilustradora) e explique o papel de cada uma.
Por que se pode falar em \textit{autoria coletiva}?
\item Segundo a nota, ``Escutar as histórias dos anciões, gravá-las,
escrevê-las e reproduzi-las permite que crianças e adultos aprendam a
ler e escrever na língua hup''. Por que é importante que uma escola
indígena ensine a ler e a escrever na língua do próprio povo?
\item Liste três coisas que se ganham e três que se perdem quando uma
narrativa oral se torna um livro impresso. Existe alguma tecnologia que
ajude a recuperar parte do que se perde? Qual?
\item Observe o livro como objeto: o texto em hup e o texto em português
aparecem juntos, na mesma página. Que mensagem essa escolha editorial
transmite ao leitor?
\end{enumerate}

\section{outras gentes: o mundo segundo os hupd'äh}

Para muitos povos indígenas, a fronteira entre ``humanos'' e ``não
humanos'' não é traçada do mesmo modo que na cultura ocidental. Na
história que você leu, a mata é habitada por outras gentes: a
gente-sombra, a gente-onça, a gente-árvore. Elas têm roupas, caçam,
conhecem cantos e benzimentos. Por isso, diz a nota, ``é preciso ter
cuidado e respeitar essas outras gentes, para que elas não façam mal aos
viajantes''.

Pesquisadores que conviveram com os Hupd'äh descrevem um universo
organizado em camadas, em que luz e sombra são forças presentes em todos
os seres, em proporções diferentes. A floresta, nessa visão, não é um
cenário vazio nem um simples estoque de recursos: é um espaço de
relações, onde é preciso saber como se comportar, quem encontrar, o que
comer, por onde andar. Nesse sentido, as narrativas funcionam também
como um saber prático, transmitido de geração em geração.

É importante lembrar que essas explicações são \textit{conhecimentos} de
um povo sobre o mundo, e não ``crendices''. Assim como outras tradições
culturais e religiosas, elas orientam condutas, valores e modos de viver
coletivamente.

\begin{enumerate}[label=\arabic*.,resume]
\item Na história, o homem-sombra canta, colhe frutas, persegue, se
confunde e é mordido. Em que ele se parece com um ser humano? Em que se
diferencia?
\item Que atitudes em relação à mata e às ``outras gentes'' a narrativa
e a nota ensinam? Você vê relação entre essas atitudes e o cuidado com o
ambiente?
\item Em muitas famílias brasileiras, circulam histórias de seres que
habitam matas, rios e estradas (a caipora, a mãe-d'água, o
lobisomem\ldots). Por que tantas culturas contam histórias sobre seres
que vivem em lugares pouco controlados pelos humanos? Que funções essas
histórias podem ter?
\end{enumerate}

\section{muitas línguas, muitos mundos}

Segundo o Censo do IBGE de 2010, falam-se no Brasil 274 línguas
indígenas, pertencentes a dezenas de famílias linguísticas. Cada língua
guarda uma forma própria de nomear e organizar o mundo, além de
narrativas, cantos, conhecimentos sobre plantas, animais e territórios.
Quando uma língua deixa de ser falada, tudo isso corre o risco de se
perder. Por essa razão, a Unesco declarou o período de 2022 a 2032 como a
Década Internacional das Línguas Indígenas.

O Alto Rio Negro, onde vivem os Hupd'äh, é uma das regiões de maior
diversidade linguística do planeta. No município de São Gabriel da
Cachoeira (AM), convivem povos de várias famílias linguísticas, e uma
lei municipal de 2002 tornou o nheengatu, o tukano e o baniwa línguas
cooficiais, ao lado do português. A língua hup pertence à família
naduhup, junto com as línguas dos Yuhupdëh, dos Dâw e dos Nadëb. Por
muito tempo, esses povos foram chamados de ``Maku'', um nome dado por
outros grupos da região e considerado pejorativo. O modo como um povo é
nomeado pelos outros também é uma questão de respeito.

\begin{enumerate}[label=\arabic*.,resume]
\item Por que a preservação de uma língua indígena interessa não só ao
povo que a fala, mas a toda a humanidade?
\item O que significa uma língua ser ``cooficial'' em um município? Que
direitos isso pode garantir aos falantes (em escolas, hospitais,
repartições públicas)?
\item Pesquise: além de São Gabriel da Cachoeira, outros municípios
brasileiros já cooficializaram línguas indígenas? Quais?
\item Por que é importante chamar um povo pelo nome que ele usa para si
(como Hupd'äh) e não por apelidos dados por outros?
\end{enumerate}

\section{ver além dos estereótipos}

A história de Way Naku mostra um povo que caça, pesca, canta, dança,
viaja, faz festas, cria escolas, escreve dicionários e publica livros. No
entanto, muitas pessoas ainda imaginam os indígenas como povos ``do
passado'', todos iguais, vivendo apenas no meio da floresta, ou como se
deixassem de ser indígenas ao usar celular, estudar na universidade ou
morar na cidade. Essas ideias são estereótipos, isto é, imagens
simplificadas e repetidas que apagam a diversidade real das pessoas e
alimentam preconceitos.

O Censo do IBGE de 2022 contou 1.694.836 pessoas indígenas no Brasil,
vivendo em todas as unidades da federação, em aldeias e cidades. A
Constituição Federal de 1988 reconhece aos povos indígenas ``sua
organização social, costumes, línguas, crenças e tradições'' (art.~231),
e a Lei 11.645/2008 tornou obrigatório o ensino da história e das
culturas indígenas nas escolas.

\begin{enumerate}[label=\arabic*.,resume]
\item Volte à lista de palavras que a turma fez antes da leitura. Quais
delas a leitura do livro e desta seção ajudaram a questionar? Quais você
acrescentaria agora?
\item Explique, com suas palavras, por que a frase ``índio de verdade
não usa celular'' é preconceituosa.
\item Um povo que registra suas histórias em livros, com a ajuda de
linguistas e antropólogos, deixa de ser tradicional? Argumente.
\end{enumerate}

\chapter{Criação e protagonismo}

Nesta etapa, você deixa de ser apenas leitor e passa a dialogar com o
livro por meio da criação. Antes de começar, lembre-se de dois cuidados
que valem para todas as propostas:

\begin{itemize}
\item a história pertence ao povo Hupd'äh: ao recriá-la, deixe claro que
se trata de uma obra inspirada em uma narrativa tradicional, citando
sempre a fonte;
\item cantos, pinturas e grafismos indígenas têm significados próprios e,
muitas vezes, sagrados. Não os copie como simples enfeite. Crie a partir
do que você entendeu, com respeito.
\end{itemize}

\section{1. novas ilustrações}

As ilustrações de Anita Ekman são uma leitura visual da história. Agora
é a sua vez.

\begin{enumerate}[label=\arabic*.]
\item Escolha um momento da narrativa que, na sua opinião, merece uma
nova imagem (por exemplo, a noite inteira em que o homem hup escuta os
cantos, ou o instante em que o homem-sombra cai da árvore).
\item Antes de desenhar, pesquise imagens reais da cucura, do umari, do
jacaré e de um igarapé do Alto Rio Negro, para que a paisagem seja fiel
ao ambiente.
\item Decida como representar a gente-sombra. A história diz que ela usa
``várias roupas de cores diferentes'' e que uma delas ``tem cor de
sombra''. Como mostrar algo que se esconde?
\item Use a técnica que preferir: lápis, colagem, pintura, desenho
digital.
\item Escreva uma legenda de três a cinco linhas explicando suas
escolhas e indicando o trecho do livro que a imagem ilustra.
\item Organizem, com a turma, uma pequena exposição das ilustrações.
\end{enumerate}

\section{2. outro ponto de vista, outro final}

\begin{enumerate}[label=\arabic*.]
\item Reconte a história do ponto de vista de Way Naku ou do jacaré, em
primeira pessoa. O que o homem-sombra pensou ao perceber que alguém o
escutava? O que sentiu ao cair da árvore?
\item Ou, se preferir, crie uma continuação: o que acontece quando o
homem hup canta os \textit{caapivaiá} pela primeira vez na aldeia? Way
Naku volta a aparecer?
\item Mantenha elementos essenciais da narrativa original: a mata, o
igarapé, a cucura, os cantos, a transmissão do conhecimento.
\item Use recursos da oralidade que você observou no livro: repetições,
onomatopeias, frases curtas, fórmulas como ``Dizem que\ldots'' ou ``Foi
isso o que ouvi\ldots''.
\item Leia seu texto em voz alta para um grupo de colegas e ouça as
sugestões antes de fazer a versão final.
\end{enumerate}

\section{3. entrevista com um personagem}

Em trios, preparem e encenem uma entrevista com o homem hup, anos depois
do encontro, quando ele já ensinou os cantos a toda a aldeia.

\begin{enumerate}[label=\arabic*.]
\item Um estudante faz o papel do entrevistador, outro o do homem hup e
o terceiro registra (em vídeo, áudio ou por escrito).
\item Elaborem pelo menos oito perguntas. Exemplos: ``O que o senhor
sentiu quando ouviu os cantos pela primeira vez?'', ``Por que bateu com
o terçado na árvore?'', ``Como ensinou os cantos aos seus filhos e
cunhados?'', ``O senhor se arrepende de algo?''.
\item As respostas devem ser coerentes com a história e com as
informações sobre os Hupd'äh que vocês estudaram.
\item Apresentem a entrevista para a turma e discutam as diferentes
interpretações que surgiram.
\end{enumerate}

\section{4. pesquisa complementar}

Em grupos, escolham um dos temas e preparem uma apresentação de até dez
minutos, com fontes identificadas.

\begin{enumerate}[label=\alph*)]
\item Os Hupd'äh hoje: onde vivem, como estão organizadas as escolas das
aldeias, quais são seus principais desafios (saúde, território, acesso a
documentos e serviços na cidade de São Gabriel da Cachoeira).
\item A família linguística naduhup e as línguas yuhup, dâw e nadëb.
\item A floresta do livro: a cucura (\textit{Pourouma cecropiaefolia}),
o umari, a cutia, o jacaré e os igarapés. Que relações ecológicas unem
esses seres?
\item Músicas e festas indígenas no Brasil: escolham dois povos e
pesquisem o papel dos cantos em suas celebrações.
\end{enumerate}

{\raggedright Consultem fontes confiáveis, como a enciclopédia Povos
Indígenas no Brasil (\url{https://pib.socioambiental.org}), publicações
de universidades e materiais produzidos por organizações indígenas.
Comparem as informações de pelo menos duas fontes.\par}

\section{5. roda de cantos e histórias}

A história de Way Naku termina afirmando que os cantos passaram ``para
seus filhos, irmãos, cunhados''. Também nas nossas famílias existem
cantos, histórias e saberes que passam de geração em geração.

\begin{enumerate}[label=\arabic*.]
\item Converse com um familiar ou vizinho mais velho e peça que conte
uma história antiga ou ensine uma cantiga que aprendeu na infância.
\item Pergunte de quem essa pessoa ouviu a história ou a cantiga, onde e
quando. Peça autorização para gravar ou anotar.
\item Transcreva a história ou a letra, mantendo o jeito de falar de
quem contou.
\item Em sala, organizem uma roda de escuta ou um pequeno podcast em que
cada estudante apresenta o que recolheu, sempre informando quem contou.
\end{enumerate}

\chapter{Pesquisa e debate}

\section{mesa-redonda: registrar é preservar?}

A turma vai simular uma mesa-redonda sobre a pergunta: \textit{Escrever
e publicar uma narrativa oral é uma forma de preservá-la ou de
transformá-la?}

\begin{enumerate}[label=\arabic*.]
\item Dividam-se em cinco grupos. Cada grupo representa uma voz
envolvida em livros como este: (1) um ancião narrador; (2) um professor
de uma escola indígena; (3) uma linguista; (4) uma editora que publica o
livro para leitores de todo o país; (5) jovens indígenas que usam
celulares e redes sociais para gravar e divulgar cantos.
\item Cada grupo prepara, com base no livro e em pesquisa, uma fala de
três minutos defendendo seu ponto de vista: quais são os ganhos, os
riscos e os cuidados necessários?
\item Um estudante atua como mediador, controla o tempo e garante que
todos falem e sejam ouvidos.
\item Após as falas, abre-se uma rodada de perguntas entre os grupos.
\item Ao final, cada estudante escreve um parágrafo com sua posição
pessoal, usando pelo menos um argumento ouvido no debate.
\end{enumerate}

\section{proposta de redação}

A partir da leitura dos textos motivadores e com base nos conhecimentos
construídos ao longo de sua formação, redija um texto
dissertativo-argumentativo, em modalidade escrita formal da língua
portuguesa, sobre o tema \textbf{``Os desafios para a preservação das
línguas indígenas no Brasil''}. Apresente proposta de intervenção que
respeite os direitos humanos. Selecione, organize e relacione, de forma
coerente e coesa, argumentos e fatos para a defesa de seu ponto de vista.

{\formular\textbf{Texto I}}
\begin{quote}
A língua hup é a primeira a ser falada pelas crianças Hupd'äh. Trata-se
de uma língua tonal, muito diferente do português. [\ldots] Escutar as
histórias dos anciões, gravá-las, escrevê-las e reproduzi-las permite
que crianças e adultos aprendam a ler e escrever na língua hup.

\hfill ``Nota dos organizadores'', \textit{Os cantos do homem-sombra}.
\end{quote}

{\formular\textbf{Texto II}}
\begin{quote}
O Censo do IBGE de 2010 identificou 305 povos indígenas e 274 línguas
indígenas faladas no Brasil. Muitas delas contam hoje com poucos
falantes.
\end{quote}

{\formular\textbf{Texto III}}
\begin{quote}
A Unesco proclamou o período de 2022 a 2032 como a Década Internacional
das Línguas Indígenas, com o objetivo de chamar a atenção para a perda
de línguas indígenas em todo o mundo e para a necessidade de
preservá-las, revitalizá-las e promovê-las.
\end{quote}

{\formular\textbf{Texto IV}}
\begin{quote}
O ensino fundamental regular será ministrado em língua portuguesa,
assegurada às comunidades indígenas também a utilização de suas línguas
maternas e processos próprios de aprendizagem.

\hfill Constituição Federal de 1988, art.~210, §~2º.
\end{quote}

Vale lembrar que a redação do Enem de 2022 teve como tema ``Desafios
para a valorização de comunidades e povos tradicionais no Brasil''.
Temas ligados à diversidade cultural e aos direitos dos povos indígenas
fazem parte do repertório esperado de quem conclui o Ensino Médio.

\chapter{Conexões com Goiás}

Os Hupd'äh vivem longe de Goiás, mas as questões que o livro levanta
--- cantos, festas, línguas, relação com a natureza, preconceito --- estão
presentes também no nosso estado. Escolha uma ou mais atividades.

\begin{enumerate}[label=\arabic*.]
\item {\formular\textbf{Povos indígenas de Goiás hoje.}} Pesquise quais
povos indígenas vivem atualmente em Goiás e onde estão suas terras.
Comece pelos Iny Karajá, que vivem às margens do rio Araguaia, em
Aruanã; pelos Avá-Canoeiro, na região de Minaçu e Colinas do Sul; e
pelos Tapuia, da Terra Indígena Carretão, entre Rubiataba e Nova
América. Para cada povo, registre: língua (ou línguas) falada(s),
localização, população aproximada e uma manifestação cultural
importante. Indique as fontes consultadas.
\item {\formular\textbf{Cantos e festas.}} Assim como os
\textit{caapivaiá} dos Hupd'äh, as festas e os rituais dos Iny Karajá
envolvem cantos e danças, como a festa do Hetohoky. Pesquise o que é
essa celebração, quem participa e qual é o seu significado. Compare com
o que você aprendeu sobre os cantos hupd'äh: em que se aproximam e em
que se diferenciam?
\item {\formular\textbf{Patrimônio.}} Os saberes e modos de fazer as
bonecas de cerâmica karajá, as \textit{ritxoko}, foram registrados pelo
Iphan como patrimônio cultural imaterial do Brasil em 2012. Pesquise o
que essas bonecas representam e por que o registro como patrimônio é
importante para o povo Karajá.
\item {\formular\textbf{Palavras indígenas no mapa de Goiás.}} Muitos
nomes de cidades, rios e lugares de Goiás têm origem indígena: o próprio
nome do estado lembra o povo Goyá, e Caiapônia lembra os Kayapó que
viviam no sul do território. Faça uma lista de dez topônimos goianos de
possível origem indígena (por exemplo, Araguaia, Jataí, Itumbiara,
Uruaçu) e pesquise o significado de cada um. O que esses nomes revelam
sobre a presença indígena no território antes e depois da colonização?
\item {\formular\textbf{Histórias da mata e do Cerrado.}} Pergunte a
pessoas mais velhas da sua família ou comunidade se conhecem histórias
de seres que vivem nas matas, nos rios ou no Cerrado goiano. Registre
uma dessas histórias e compare-a com a narrativa de Way Naku.
\end{enumerate}

\chapter{Projeto integrador}

\section{vozes que ensinam: cantos, histórias e memórias}

{\formular\textbf{Objetivo.}} Produzir, em grupo, uma mostra sonora e
visual que reúna narrativas e cantos indígenas estudados, histórias e
cantigas recolhidas na comunidade e criações da turma, valorizando a
oralidade como forma de conhecimento.

{\formular\textbf{Componentes curriculares envolvidos.}} Língua
Portuguesa, Arte, História, Geografia, Sociologia e Biologia.

{\formular\textbf{Etapas.}}

\begin{enumerate}[label=\arabic*.]
\item {\formular\textbf{Planejamento}} (Língua Portuguesa e Sociologia):
a turma define o formato final (exposição, sarau, série de podcasts ou
combinação desses formatos), divide tarefas e combina prazos.
\item {\formular\textbf{Pesquisa}} (História e Geografia): os grupos
reúnem as pesquisas feitas nas atividades anteriores sobre os Hupd'äh,
sobre a diversidade linguística do Alto Rio Negro e sobre os povos
indígenas de Goiás, elaborando mapas e linhas do tempo.
\item {\formular\textbf{Natureza e cultura}} (Biologia): um grupo prepara
painéis sobre as plantas e os animais citados no livro e sobre a relação
entre os saberes indígenas e a conservação da floresta.
\item {\formular\textbf{Criação}} (Arte): seleção das ilustrações, dos
recontos e das entrevistas produzidos; ensaio de leituras em voz alta e
de cantigas recolhidas com as famílias; criação de uma trilha sonora
com percussão corporal ou instrumentos simples, inspirada no ritmo dos
cantos de abertura do livro (sem imitar os cantos hupd'äh).
\item {\formular\textbf{Montagem e revisão}}: produção de legendas,
créditos e fichas técnicas, com identificação de todas as fontes e de
todas as pessoas que contaram histórias.
\item {\formular\textbf{Socialização}}: apresentação para outras turmas
e para as famílias, seguida de uma roda de conversa.
\end{enumerate}

{\formular\textbf{Autoavaliação.}} Ao final, responda no diário de
leitura:

\begin{itemize}
\item Participei das etapas combinadas e cumpri os prazos?
\item Minhas produções respeitaram a autoria e o significado das
narrativas indígenas?
\item O que eu pensava sobre os povos indígenas antes deste estudo, e o
que penso agora?
\item Que pergunta sobre os Hupd'äh ou sobre os povos indígenas de Goiás
eu ainda gostaria de investigar?
\end{itemize}

\chapter{Para saber mais}
\begingroup\raggedright

\section{livros}

\begin{itemize}
\item \textsc{krenak}, Ailton. \textit{Ideias para adiar o fim do
mundo}. São Paulo: Companhia das Letras, 2019. Reflexões de um
pensador indígena sobre a relação entre humanidade e natureza.
\item \textsc{munduruku}, Daniel. \textit{Contos indígenas brasileiros}.
São Paulo: Global, 2004. Narrativas de diferentes povos, reunidas por um
escritor indígena.
\item \textsc{ramos}, Danilo Paiva. \textit{Círculos de coca e fumaça}.
São Paulo: Hedra (Coleção Mundo Indígena). Estudo sobre os encontros
noturnos e os caminhos dos Hupd'äh, escrito por um dos organizadores
deste livro.
\item \textsc{jecupé}, Kaká Werá. \textit{A terra dos mil povos}. São
Paulo: Peirópolis, 1998. Uma história do Brasil contada do ponto de
vista indígena.
\item \textit{A árvore dos cantos} e \textit{A terra uma só}, da Coleção
Mundo Indígena, que você lerá nas próximas séries.
\end{itemize}

\section{filmes}

\begin{itemize}
\item \textit{A febre} (Maya Da-Rin, 2019). Ficção protagonizada por
indígenas do Alto Rio Negro, sobre um trabalhador desana que vive em
Manaus.
\item Filmes do projeto Vídeo nas Aldeias
(\url{https://videonasaldeias.org.br}), realizados por cineastas
indígenas de diversos povos.
\end{itemize}

\section{sites}

\begin{itemize}
\item Povos Indígenas no Brasil, do Instituto Socioambiental:
\url{https://pib.socioambiental.org}.
\item Terras Indígenas no Brasil: \url{https://terrasindigenas.org.br}.
\item Articulação dos Povos Indígenas do Brasil (Apib):
\url{https://apiboficial.org}.
\end{itemize}
\endgroup
